\chapter{System Design}

\section{System Architecture}
BHUVAN follows a \textbf{client-only single-page application (SPA)} architecture. There is no backend server---all computation (terrain generation, physics simulation, hazard analysis, audio synthesis) executes in the user's browser using JavaScript and WebGL.

\begin{figure}[h]
\centering
\begin{tikzpicture}[node distance=1.2cm, auto,
    block/.style={rectangle, draw, fill=blue!10, text width=3cm, text centered, rounded corners, minimum height=1cm},
    engine/.style={rectangle, draw, fill=orange!15, text width=2.8cm, text centered, rounded corners, minimum height=1cm},
    arrow/.style={->, >=stealth, thick}]

    \node[block] (app) {App.jsx\\State Machine};
    \node[block, below left=1.5cm and 0.5cm of app] (intro) {IntroScreen};
    \node[block, below=1.5cm of app] (scene) {SceneCanvas\\(R3F Canvas)};
    \node[block, below right=1.5cm and 0.5cm of app] (hud) {HUD Overlay};

    \node[engine, below left=1.5cm and 0cm of scene] (terrain_e) {terrain.js\\Heightmap Gen};
    \node[engine, below=1.5cm of scene] (physics_e) {physics.js\\Lander Dynamics};
    \node[engine, below right=1.5cm and 0cm of scene] (audio_e) {audio.js\\Web Audio};

    \node[block, below left=1.2cm and -0.5cm of terrain_e] (terrain_c) {Terrain.jsx\\Mesh + Colors};
    \node[block, below=1.2cm of physics_e] (lander_c) {Lander.jsx\\Model + Particles};

    \draw[arrow] (app) -- (intro);
    \draw[arrow] (app) -- (scene);
    \draw[arrow] (app) -- (hud);
    \draw[arrow] (scene) -- (terrain_e);
    \draw[arrow] (scene) -- (physics_e);
    \draw[arrow] (app) -- (audio_e);
    \draw[arrow] (terrain_e) -- (terrain_c);
    \draw[arrow] (physics_e) -- (lander_c);
\end{tikzpicture}
\caption{System Architecture Diagram}
\label{fig:architecture}
\end{figure}

\section{Application State Machine}
The application progresses through four distinct phases:

\begin{figure}[h]
\centering
\begin{tikzpicture}[node distance=3cm, auto,
    state/.style={circle, draw, fill=green!10, minimum size=1.8cm, text centered, font=\small},
    arrow/.style={->, >=stealth, thick}]

    \node[state] (intro) {INTRO};
    \node[state, right of=intro] (orbital) {ORBITAL};
    \node[state, right of=orbital] (descent) {DESCENT};
    \node[state, above right=0.8cm and 2cm of descent] (landed) {LANDED};
    \node[state, below right=0.8cm and 2cm of descent] (crashed) {CRASHED};

    \draw[arrow] (intro) -- node[above]{\small Begin} (orbital);
    \draw[arrow] (orbital) -- node[above]{\small Initiate} (descent);
    \draw[arrow] (descent) -- node[above,sloped]{\small $|v_y|<3$} (landed);
    \draw[arrow] (descent) -- node[below,sloped]{\small $|v_y|\geq3$} (crashed);
\end{tikzpicture}
\caption{Application State Machine}
\label{fig:state_machine}
\end{figure}

\section{Use Case Diagram}

\begin{figure}[h]
\centering
\begin{tikzpicture}
    \node[draw, minimum width=9cm, minimum height=7cm, label=above:{\textbf{BHUVAN System}}] (sys) at (5,3.5) {};
    
    \node[draw, ellipse, text width=2.5cm, text centered] (uc1) at (5,6.2) {View Intro};
    \node[draw, ellipse, text width=2.5cm, text centered] (uc2) at (3,4.8) {Survey Terrain};
    \node[draw, ellipse, text width=2.5cm, text centered] (uc3) at (7,4.8) {Switch View Mode};
    \node[draw, ellipse, text width=2.5cm, text centered] (uc4) at (3,3.2) {Select Landing Zone};
    \node[draw, ellipse, text width=2.5cm, text centered] (uc5) at (7,3.2) {Toggle Autopilot};
    \node[draw, ellipse, text width=2.5cm, text centered] (uc6) at (5,1.8) {Execute Descent};
    \node[draw, ellipse, text width=2.5cm, text centered] (uc7) at (5,0.5) {View Results};

    \node (user) at (-1.5,3.5) {\includegraphics[width=0.8cm]{example-image}};
    \node at (-1.5,2.8) {User};

    \draw (user) -- (uc1);
    \draw (user) -- (uc2);
    \draw (user) -- (uc3);
    \draw (user) -- (uc4);
    \draw (user) -- (uc5);
    \draw (user) -- (uc6);
    \draw (user) -- (uc7);
\end{tikzpicture}
\caption{Use Case Diagram}
\label{fig:usecase}
\end{figure}

\section{Component Design}

\begin{table}[h]
\centering
\caption{Component Responsibilities}
\label{tab:components}
\begin{tabular}{|p{3.5cm}|p{3cm}|p{6cm}|}
\hline
\textbf{Component} & \textbf{Type} & \textbf{Responsibility} \\
\hline
App.jsx & React Component & State machine, keyboard input, physics loop orchestration \\
\hline
IntroScreen.jsx & React Component & CRT boot animation, mission briefing display \\
\hline
HUD.jsx & React Component & Telemetry gauges, bars, warnings, view mode controls \\
\hline
SceneCanvas.jsx & R3F Canvas & 3D scene composition, camera rig, post-processing \\
\hline
Terrain.jsx & R3F Mesh & Heightmap mesh with switchable vertex colors \\
\hline
Lander.jsx & R3F Group & Geometric lander model, thruster flame, exhaust particles \\
\hline
terrain.js & Pure Module & fBm noise, crater generation, slope computation, sampling \\
\hline
physics.js & Pure Module & Newtonian dynamics, autopilot controller, collision detection \\
\hline
audio.js & Pure Module & Web Audio oscillators, noise generators, gain control \\
\hline
\end{tabular}
\end{table}

\section{User Interface Design}
The UI follows a \textbf{mission-control aesthetic} inspired by SpaceX and NASA flight displays:
\begin{itemize}
    \item \textbf{Color palette:} Dark background (\texttt{\#0a0a0f}), green accents (\texttt{\#00ff88}), orange warnings (\texttt{\#ff8800}), red critical (\texttt{\#ff3344}).
    \item \textbf{Typography:} JetBrains Mono (monospace, telemetry), Outfit (display headings).
    \item \textbf{Layout:} Full-screen 3D canvas with glassmorphic HUD panels overlaid at edges.
    \item \textbf{Animations:} CSS keyframe animations for warning pulses, button breathing, and CRT scanlines.
\end{itemize}

\section{Performance Considerations}
\begin{itemize}
    \item Terrain mesh uses a single \texttt{PlaneGeometry} with 65,536 vertices (256$\times$256)---well within WebGL vertex limits.
    \item Vertex colors are computed once at load time; view mode changes trigger a single geometry rebuild, not per-frame computation.
    \item Post-processing uses mipmapped bloom to reduce shader passes.
    \item Particle system uses a fixed pool of 80 particles with manual buffer updates---no object allocation per frame.
    \item OrbitControls uses damping (\texttt{dampingFactor=0.08}) to reduce input event frequency.
\end{itemize}
